\documentclass[aspectratio=169]{beamer}

\usetheme{Madrid}
\usecolortheme{default}

\usepackage[T1]{fontenc}
\usepackage[utf8]{inputenc}
\usepackage[brazil]{babel}
\usepackage{amsmath}
\usepackage{booktabs}
\usepackage{array}
\usepackage{tikz}
\usetikzlibrary{positioning,arrows.meta}
\usepackage{qrcode}

\title[Revisão Sistemática]{Etapas Iniciais da Revisão Sistemática da Literatura}
\subtitle{Segmentação 3D de Órgãos em TC com Deep Learning\\DOUTORAMENTO EM INFORMÁTICA (UTAD)}
\author{Alexandre T. Nogueira}
\date{}

\begin{document}

\begin{frame}
\titlepage
\vspace{-1.2em}
\begin{center}
\qrcode[hyperlink,height=2cm]{https://github.com/gameguild-gg/docdo-paper}
\end{center}
\end{frame}

\begin{frame}{Roteiro}
	\tableofcontents[hideallsubsections]
\end{frame}

\section{Resumo}
\begin{frame}{1. Resumo}
\small
\textbf{Objetivo:} identificar quais algoritmos de segmentação 3D de órgãos em Tomografia Computadorizada (TC) estão prontos para adoção clínica em planejamento cirúrgico e quais barreiras permanecem.

\vspace{0.4em}
\textbf{Método:}
\begin{itemize}
	\item Busca em 3 bases (PubMed, arXiv, Semantic Scholar), 2015--jan/2026.
	\item PRISMA modificado com triagem semi-automática (Elasticsearch + LLMs).
	\item Síntese narrativa de 52 estudos incluídos.
	\item Meta-análise formal não aplicada (heterogeneidade de protocolos).
\end{itemize}

\textbf{Principais achados:}
\begin{itemize}
	\item nnU-Net (82,3\% Dice) e Swin UNETR (82,4\%) com desempenho de ponta semelhante.
	\item CNNs dominam (98,1\%); variantes 3D U-Net em 82,7\% dos trabalhos.
	\item Gargalos de translação: \textit{domain shift}, segmentação vascular inferior e integração em fluxo clínico.
\end{itemize}
\end{frame}

\section{Introdução e Definições}
\begin{frame}{2. Introdução ao tema e definições}
\small
\textbf{Contexto clínico:}
\begin{itemize}
	\item Segmentação de órgãos em Tomografia Computadorizada (TC) é central para planejamento cirúrgico, navegação e quantificação.
	\item Motivação clínica: o contorno manual continua a demandar muito tempo e limita a adoção rotineira do planejamento cirúrgico.
\end{itemize}

\textbf{Definições-chave:}
\begin{itemize}
	\item \textbf{Segmentação semântica multi-classe}: rótulo por voxel para múltiplos órgãos.
	\item \textbf{Segmentação por instância}: separa instâncias da mesma classe, por exemplo, rim esquerdo e rim direito.
	\item \textbf{Segmentação panóptica}: combina semântica + instância.
	\item \textbf{Segmentação volumétrica (3D)}: usa contexto inter-slices (alvo principal desta revisão).
\end{itemize}
\end{frame}

\section{Questões de Investigação}
\begin{frame}{3. Questões de investigação (RQs)}
\small
\begin{block}{RQ1}
Quais paradigmas arquiteturais dominam a segmentação 3D de órgãos em TC e como diferem em princípios, vieses indutivos e custo computacional?
\end{block}

\begin{block}{RQ2}
Quais datasets e métricas são padrão para benchmark e quais vieses/limitações de generalização eles introduzem?
\end{block}

\begin{block}{RQ3}
Quais \textit{trade-offs} existem entre acurácia, custo computacional e esforço de anotação? Como está a generalização entre domínios?
\end{block}

\begin{block}{RQ4}
Quais lacunas impedem a adoção clínica (reprodutibilidade, \textit{domain shift}, integração e barreiras regulatórias)?
\end{block}
\end{frame}

\section{Palavras-chave}
\begin{frame}{4. Palavras-chave}
\small
\begin{itemize}
	\item medical image segmentation
	\item computed tomography
	\item deep learning
	\item convolutional neural networks
	\item vision transformers
	\item 3D volumetric segmentation
	\item organ segmentation
	\item surgical planning
	\item clinical deployment
\end{itemize}
\end{frame}

\section{Base de Dados}
\begin{frame}{5. Base de dados para pesquisa}
\small
\begin{table}
\centering
\begin{tabular}{>{\raggedright\arraybackslash}p{2.8cm} >{\raggedright\arraybackslash}p{5.4cm} >{\raggedright\arraybackslash}p{2.4cm}}
\toprule
\textbf{Base} & \textbf{Acesso} & \textbf{Registros brutos} \\
\midrule
PubMed & Entrez E-utilities (API oficial NCBI) & 997 \\
arXiv & API oficial arXiv & 904 \\
Semantic Scholar & Graph API v1 & 1.084 \\
\midrule
\textbf{Total} & \textbf{Antes de deduplicação} & \textbf{2.985} \\
\bottomrule
\end{tabular}
\end{table}

\vspace{0.3em}
\footnotesize
Escopo confirmado pelos artefatos suplementares: deduplicação para 2.821 estudos únicos.
\end{frame}

\section{String de Pesquisa}
\begin{frame}{6. String de pesquisa}
\small
\textbf{Construção por 3 conceitos:}
\begin{itemize}
	\item \textbf{Imagem/população}: \texttt{``CT'' OR ``computed tomography'' OR ``CT scan'' OR ``CT imaging''}
	\item \textbf{Tarefa}: \texttt{``segmentation'' AND (``3D'' OR ``volumetric'' OR ``organ'' OR ``multi-organ'' OR ``abdominal'')}
	\item \textbf{Método}: \texttt{``deep learning'' OR ``neural network'' OR ``convolutional'' OR ``U-Net'' OR ``transformer'' OR ``encoder-decoder''}
\end{itemize}

\textbf{Exemplo (PubMed, resumido):}
\begin{block}{Consulta Booleana}
\scriptsize
\begin{itemize}
	\item \texttt{(computed tomography OR CT)}
	\item \texttt{AND (image segmentation OR segmentation)}
	\item \texttt{AND (3D OR volumetric OR organ OR multi-organ OR abdominal)}
	\item \texttt{AND (deep learning OR neural network* OR convolutional OR U-Net OR transformer OR encoder-decoder)}
	\item \texttt{AND 2015:2026[dp] AND English[la]}
\end{itemize}
\end{block}
\end{frame}

\section{Critérios de Inclusão e Exclusão}
\begin{frame}{7. Critérios de inclusão e exclusão}
\scriptsize
\begin{columns}[T]
\column{0.49\textwidth}
\textbf{Inclusão --- todos obrigatórios:}
\begin{itemize}
	\item[\textbf{1}] DL: CNN, Transformer, U-Net, encoder-decoder
	\item[\textbf{2}] 3D/volumétrico: conv.\ 3D ou 2.5D c/ contexto inter-slice
	\item[\textbf{3}] TC incluída (mono ou multi-modal c/ CT)
	\item[\textbf{4}] Segmentação de órgãos anatômicos (fígado, rim, pâncreas, pulmão\ldots)
	\item[\textbf{5}] Contribuição metodológica original
	\item[\textbf{6}] 2015--jan/2026; inglês; métrica quantitativa (Dice, HD\ldots)
\end{itemize}

\column{0.49\textwidth}
\textbf{Exclusão --- qualquer um exclui:}
\begin{itemize}
	\item[\textbf{1}] Só detecção/classificação (sem segmentação espacial)
	\item[\textbf{2}] Modalidade não-TC exclusiva (MRI, US, raio-X\ldots)
	\item[\textbf{3}] Alvo só tumor/vaso/osso sem órgão parênquima
	\item[\textbf{4}] Método 2D puro sem contexto volumétrico
	\item[\textbf{5}] ML clássico / sem DL
	\item[\textbf{6}] Revisão/survey sem método novo
\end{itemize}
\end{columns}
\end{frame}

\section{Diagrama de Fluxo PRISMA}
\begin{frame}{8. Diagrama de fluxo (PRISMA)}
\centering
\scriptsize
\begin{tikzpicture}[
		node distance=0.45cm,
		box/.style={rectangle, draw, rounded corners, align=center, minimum width=8.3cm, minimum height=0.55cm, inner sep=2.5pt},
		side/.style={rectangle, draw, align=center, minimum width=3.6cm, minimum height=0.55cm, inner sep=2pt},
		arr/.style={-{Stealth[length=2mm]}, thick}
]
	\node[box] (s1raw) {Registros identificados (PubMed + arXiv + Semantic Scholar): \textbf{2.985}};
	\node[box, below=of s1raw] (dedup) {Após deduplicação (DOI + título normalizado): \textbf{2.821}};
	\node[side, right=0.2cm of dedup] (dedupexc) {Duplicatas removidas: \textbf{164}};
	\node[box, below=of dedup] (s2) {Pré-filtro booleano Elasticsearch (titulo e resumo): \textbf{638}};
	\node[side, right=0.2cm of s2] (s2exc) {Excluídos no pré-filtro: \textbf{2.183}\\
	\tiny must: DL; segm.; CT; órgão; 3D/vol.\\
	\tiny excl.: revisão/survey no título};
	\node[box, below=of s2] (llm) {Triagem via LLM (titulo e resumo): \textbf{161}};
	\node[side, right=0.2cm of llm] (llmexc) {Excluídos via LLM: \textbf{477}\\
	\tiny mod.~não-CT; 2D sem vol.; não-órgão;\\
	\tiny sem DL; revisão sem método novo};
	\node[box, below=of llm] (pdf) {Texto completo obtido (PDF): \textbf{63}};
	\node[side, right=0.2cm of pdf] (nopdf) {Sem PDF:\textbf{98}};
	\node[box, below=of pdf] (s3) {Triagem de texto completo: \textbf{52 incluídos}};
	\node[side, right=0.2cm of s3] (s3exc) {Excluídos (texto completo): \textbf{11}\\
	\tiny CBCT:2; 2D sem vol.:2; lesão/tumor:2;\\
	\tiny sem Dice/HD:2; sem contrib.:3};
	\node[box, below=of s3] (final) {Estudos na síntese narrativa final: \textbf{52}};

	\draw[arr] (s1raw) -- (dedup);
	\draw[arr] (dedup) -- (s2);
	\draw[arr] (s2) -- (llm);
	\draw[arr] (llm) -- (pdf);
	\draw[arr] (pdf) -- (s3);
	\draw[arr] (s3) -- (final);
	\draw[arr] (dedup.east) -- (dedupexc.west);
	\draw[arr] (s2.east) -- (s2exc.west);
	\draw[arr] (llm.east) -- (llmexc.west);
	\draw[arr] (pdf.east) -- (nopdf.west);
	\draw[arr] (s3.east) -- (s3exc.west);
\end{tikzpicture}

\end{frame}

\begin{frame}{Próximos passos}
\small
\begin{columns}[c]
\column{0.65\textwidth}
\begin{itemize}
  \item \textbf{Validar} o pipeline de triagem automatizada em subconjunto com rótulos manuais (precisão, recall, F1).
  \item \textbf{Comparar} os modelos mais citados (TotalSegmentator, nnU-Net, MedSAM) nos datasets públicos levantados.
  \item \textbf{Integrar ao 3D Slicer e DocDo} em três fases:
  \begin{enumerate}\scriptsize
    \item \textit{Baseline}: nnU-Net/TotalSegmentator para segmentação robusta sem ajuste sítio-específico.
    \item \textit{Qualidade}: MC-Dropout/TTA para sinalizar casos de baixa confiança para revisão.
    \item \textit{Refinamento}: MedSAM com prompts de bounding-box para correção interativa pelo cirurgião.
  \end{enumerate}
\end{itemize}
\column{0.30\textwidth}
\centering
\qrcode[hyperlink,height=4.5cm]{https://github.com/gameguild-gg/docdo-paper}
\end{columns}
\end{frame}

\end{document}
